\section{Related Work}
\label{sec:related}

\noindent\textbf{The controller this paper deploys.}
Self-regulating cars learns a policy that assigns a desired speed to each
super-segment of a highway network, and reports a 5\% throughput gain and a 13\%
delay reduction against no control in PTV Vissim~\cite{self_regulating_cars}.
Its observer is central and its decisions reach vehicles through a cloud
service. The failure mode it targets is characterized
in~\cite{bhardwaj2023understanding}, which defines a sudden traffic jam as a
burst that pushes density past a critical threshold into a stable low-capacity
state. This paper changes where the policy is evaluated and what it is allowed
to read, and it leaves the reward, the training and the controller unchanged.

\noindent\textbf{Sparse vehicles as traffic actuators.}
A single controlled vehicle damps stop-and-go waves on a ring
track~\cite{stern2018dissipation}, which established the sparse-actuator
premise. Learned controllers improve simulated traffic at low penetration in the
mixed-autonomy literature~\cite{wu2017flow,vinitsky2018benchmarks,kreidieh2018dissipating}.
CIRCLES placed about one hundred controlled vehicles into live interstate
traffic and is the strongest field evidence available, using purpose-built
vehicles and a centralized speed planner carried over cellular
links~\cite{lee2024traffic}. The work here differs on the deployment axis rather
than the control axis. The vehicle is unmodified, the compute is commodity, and
no planner is contacted at run time.

\noindent\textbf{Infrastructure-centric control.}
Ramp metering, variable speed limits and adaptive signal control remain the
deployed tools~\cite{arnold1998ramp,papageorgiou2002freeway,scats,scoot}. They
are effective where the sensors and actuators exist, and they require owning the
roadside. The entry gate in Section~\ref{sec:simulation} is ramp metering
performed by the policy at the network boundary, which is the one place a
vehicle-side controller has no upstream link to slow.

\noindent\textbf{Traffic state from vehicles and from feeds.}
Probe-vehicle systems estimate network state by uplinking per-vehicle traces to
a backend that aggregates
centrally~\cite{seo2017traffic,herrera2010evaluation}. Connected-vehicle
standards broadcast per-vehicle position, speed and heading at up to ten
hertz~\cite{sae2020j2735,etsi2019cam,kenney2011dedicated}, and the collective
perception service extends that to per-object detections~\cite{etsi2023cpm}.
This design uses neither. The observation it needs is already published as a
commercial product at segment granularity~\cite{here_traffic_api}, so a query
replaces both the radio and the aggregation backend. The trade is explicit. A
feed is a third-party dependency with an unreported lag, and the two are
different risks rather than one being strictly better.

\noindent\textbf{Advisory interfaces.}
Coarse, human-followable advisories have precedent in advisory-autonomy
work~\cite{Yan_2023,cho2023temporal,fridman2018human}, and mobile telematics
platforms operate the driver-facing rail at the scale of tens of millions of
drivers~\cite{cmt2025,samsara2025}. This paper does not evaluate compliance, for
the reasons stated in Section~\ref{sec:intro}.

\noindent\textbf{Edge perception.}
Embedded detectors and Jetson-class hardware make on-vehicle perception
practical~\cite{redmon2016yolo,nvidia2023jetson,bewley2016sort,yang2024depth}.
The measurements in Section~\ref{sec:eval} are consistent with that literature.
The contribution here is not the detector. It is what the rest of the system
around it has to do to keep the detector's output honest when it is wrong, which
Section~\ref{sec:experience} reports.
